%---------------------------------------%
%-------------- Appendix ------------------%
%---------------------------------------%



\appendix
\section{Proof of the vertex and source statement in Fong's second reduction}\label{app:A}%


In this Appendix, we give a proof of Part (b) of Theorem~\ref{thm:Fong2ndReduction} establishing that Fong's second reduction can be understood as a Morita equivalence with endo-permutation source. This result, understood in terms of source-algebras, is due to Puig and we refer to \cite[Theorem~9.11.9] {LinckelmannVol2}, which states the connection between the source-algebras of the blocks involved and the source of the bimodule giving the Morita equivalence.  Below we give a precise proof exhibiting a source and a vertex of this bimodule purely in module-theoretic terms, which in particular does not require knowledge of the source-algebras. 




\begin{proof}[Proof of Theorem~\ref{thm:Fong2ndReduction}, Part {\rm(b)} and
Remark~\ref{rem:ProofFong2ndRed}]{\leavevmode}

\smallskip
\noindent
{\rm(i)} \emph{The module $V$ and its restriction to $D$.}
Since $N\cap D=1$, the quotient group $ND/N\cong D$ is a $p$-group. Moreover, since $\bc$ is $G$-stable, the inertia group of the simple $\bc$-module $V$ in $ND$ is $ND$. Therefore, by
\cite[Chapter~III, Corollary~3.16]{FeitBook}, $V$ has a unique extension, up to isomorphism, to an irreducible $k[ND]$-module. We denote this extension again by $V$ and set
\( W:=\Res_D^{ND}(V)\). 
Consider $S=kNe_{\bc}$ as a $kD$-module via conjugation. Then
\[
S=kNe_{\bc}\cong\End_k(V)\cong\End_k(W)
\]
as $kD$-modules. Since $S$ is a direct summand of the permutation $kD$-module $kN$, it follows that $\End_k(W)$ is a permutation $kD$-module. Hence $W$ is an endo-permutation $kD$-module.

Moreover, $\Br_D(e_{\bc})\neq0$, and hence $\End_k(V)\cong kNe_{\bc}$ has a trivial direct summand as a $kD$-module. It follows from \cite[Lemma~2.2]{ThevenazSurvey} that $W$ is capped. We set
\[
U:=\operatorname{cap}(W).
\]
Thus $U$ is an indecomposable $kD$-module with vertex $D$.

We next show that the $kND$-module $V$ has vertex $D$. Let $Q$ be a vertex of $V$. Since $\bc$ is a block of defect zero, the module $V$ belongs to a block of $kND$ having a defect group $P$ such that $P\cap N=1$. Therefore, $|Q|\leq |P|\leq |D|$. On the other hand, $U$ is a direct summand of $\Res_D^{ND}(V)$ and has vertex $D$. Hence $D$ is contained, up to conjugation, in a vertex of $V$. It
follows that $|D|\leq |Q|$. Therefore $|Q|=|D|$, and $D$ is a vertex of $V$.

\smallskip
\noindent
{\rm(ii)} \emph{Construction of the Morita equivalence.}
By \cite[Section~3]{SalminenSOURCES08}, there exists an isomorphism
\[
\Phi:S\otimes_k kHe\xrightarrow{\sim}kGe_{\bc}
\]
of interior $D$-algebras, where $e$ is an idempotent of $kZ$. Here, we identify the copy of $D$ in $H$ arising from this construction with $D$.

Since $S$ is simple, there is a bijection between the block idempotents of $kH$ appearing in $e$ and the blocks of $kG$ appearing in $e_{\bc}$; see
\cite[Section~4]{SalminenSOURCES08}. Let
$\widehat{\bB}$ be the block of $kH$ corresponding to
$\bB$. By the proof of
\cite[Proposition~4.1]{SalminenSOURCES08}, which does not use the nilpotency assumption, the chosen copy of $D$ in $H$ is a defect group of $\widehat{\bB}$.

The isomorphism $\Phi$ restricts to an isomorphism $\Phi_{\bB}:S\otimes_k\widehat{\bB}\xrightarrow{\sim}\bB$ of interior $D$-algebras. Since $S\cong\End_k(V)$, the $(S\otimes_k\widehat{\bB},\widehat{\bB})$-bimodule $V\otimes_k\widehat{\bB}$ induces a Morita equivalence between
$S\otimes_k\widehat{\bB}$ and $\widehat{\bB}$. Its actions are given by
\[
(s\otimes a)\cdot(v\otimes x)\cdot a'=s(v)\otimes axa'
\]
for all $s\in S$, $a,a',x\in\widehat{\bB}$ and $v\in V$; see \cite[Exercise~9.5(a)]{ThevenazBook}. Transporting the left action along $\Phi_{\bB}^{-1}$, the bimodule $M:=V\otimes_k\widehat{\bB}$ becomes a $(\bB,\widehat{\bB})$-bimodule and induces a Morita equivalence between $\bB$ and $\widehat{\bB}$.

\smallskip
\noindent
{\rm(iii)} \emph{Relative projectivity of $M$.}
Set $\overline G:=G/N$
and let $\pi:H\to\overline G$
be the canonical surjection. Consider the $(kGe_{\bc},kHe)$-bimodule $\mathcal M:=V\otimes_k kHe$, where the left action is transported along $\Phi:S\otimes_k kHe\xrightarrow{\sim}kGe_{\bc}$. 
Recall that $Z$ is identified with a finite subgroup of $k^\times$ and that
\[
e=\frac{1}{|Z|}
\sum_{\zeta\in Z}\zeta^{-1}(1,\zeta).
\]
For every $\zeta_0\in Z$, we have $(1,\zeta_0)e=\zeta_0e$, where, on the right-hand side, $\zeta_0$ is regarded as an element of $k^\times$. For every $\overline x\in\overline G$, choose an element $\widetilde x\in H$ with $\pi(\widetilde x)=\overline x$. Since $Z$ is central in $H$, we have
\[
(1,\zeta)\widetilde xe=\widetilde x(1,\zeta)e=\zeta\widetilde xe \qquad(\zeta\in Z).
\]
It follows that $kZ\widetilde xe=k\widetilde xe$. Therefore, as a $k$-vector space,
\[
kHe=\bigoplus_{\overline x\in\overline G} k\widetilde xe,
\]
and hence
\[
\mathcal M=\bigoplus_{\overline x\in\overline G}
\mathcal M_{\overline x}\;\;\text{ with }\quad\mathcal M_{\overline x}:=V\otimes_k k\widetilde xe.
\]
We now use the explicit construction of the same
isomorphism $\Phi$. By \cite[Section~3]{SalminenSOURCES08}, for every $g\in G$
there exist $s_g\in S^\times$ and $\alpha_g\in k^\times$ such that
\[
\Phi^{-1}(ge_{\bc})=s_g\otimes \alpha_g\widetilde{gN}e,
\tag{1}
\]
where $\widetilde{gN}\in H$ is a lift of
$gN\in\overline G$. 
Let $\overline x\in\overline G$, $v\otimes y\in\mathcal M_{\overline x}$ and $y\in k\widetilde xe$. Using~{\rm(1)}, we obtain
\begin{align*}
(g,h)\cdot(v\otimes y)&=g\cdot(v\otimes y)\cdot h^{-1}=\alpha_gs_g(v)\otimes \widetilde{gN}yh^{-1}.
\end{align*}
The image of
$\widetilde{gN}\widetilde xh^{-1}$ under $\pi$ is
$(gN)\overline x\pi(h)^{-1}$. Therefore $(g,h)\mathcal M_{\overline x} \subseteq \mathcal M_{(gN)\overline x\pi(h)^{-1}}$.
Since the map induced by $(g,h)$ is invertible and both spaces have dimension $\dim_k(V)$, we obtain
\[
(g,h)\mathcal M_{\overline x}=\mathcal M_{(gN)\overline x\pi(h)^{-1}}.
\tag{2}
\]
In particular, $(g,h)\mathcal M_{\overline1}=\mathcal M_{\overline1}$ if and only if $gN=\pi(h)$. Thus the stabilizer of $\mathcal M_{\overline1}=V\otimes_k ke$ in $G\times H$ is
\[
L:=\left\{(g,h)\in G\times H\ \middle|\ gN=\pi(h) \right\}.
\]
Moreover, Equation~{\rm(2)} shows that $G\times H$ acts transitively on $\left\{\mathcal M_{\overline x}\ \middle|\ \overline x\in\overline G\right\}$. It follows that the natural homomorphism
\[
\Theta:\Ind_L^{G\times H}(\mathcal M_{\overline1})\longrightarrow\mathcal M,\qquad(g,h)\otimes m\longmapsto(g,h)\cdot m,
\]
is surjective. Since both sides have the same dimension, $\Theta$ is an isomorphism, so $\mathcal M$ is induced from $\mathcal M_{\overline1}$. Since $\widehat{\bB}$ is a block summand of $kHe$, the bimodule $M$ is a direct summand of $\mathcal M$ and is therefore
relatively $L$-projective.

\smallskip
\noindent
{\rm(iv)} \emph{Vertex and source of $M$.}
Since $M$ induces a Morita equivalence between the block algebras $\bB$ and $\widehat{\bB}$, it is indecomposable as a $k[G\times H]$-module. Let $R$ be a vertex of $M$. Since $M$ is relatively $L$-projective, $R$ is contained in a conjugate of $L$. On the other hand, $M$ belongs to the block $\bB\otimes_k\widehat{\bB}^*$ of $k[G\times H]$. Hence, after replacing $R$ by a conjugate, we may assume that $R\leq (D\times D)\cap{}^yL$ for some $y\in G\times H$.

Since $D\cap N=1$ and $D\cap Z=1$, it follows from the definition of $L$ that $(D\times D)\cap{}^yL$ is a twisted diagonal
subgroup of $D\times D$.
In particular, $|R|\leq |D|$.

Viewed as a $k[H\times H]$-module, the regular bimodule $\widehat{\bB}$ has vertex $\Delta D$ and trivial source. Therefore $k\mid
\Res_{\Delta D}^{H\times H}(\widehat{\bB})$. Since $\Phi_{\bB}$ is an isomorphism of interior
$D$-algebras, the action of $\Delta D$ on $M$ is given by
\[
(d,d)\cdot(v\otimes x) =dv\otimes dxd^{-1}.
\]
Therefore,
\[
\Res_{\Delta D}^{G\times H}(M)\cong W\otimes_k
\Res_{\Delta D}^{H\times H}(\widehat{\bB})
\]
as $k[\Delta D]$-modules, where we identify
$D\cong\Delta D$. 
It follows that $W\mid\Res_{\Delta D}^{G\times H}(M)$. Since $U\mid W$, we obtain $U\mid \Res_{\Delta D}^{G\times H}(M)$.
Viewed as a $k[\Delta D]$-module, $U$ has vertex
$\Delta D$. Therefore, there exists $z\in G\times H$ such that $\Delta D\leq{}^zR$.
Hence $|D|\leq |R|$. Together with the preceding inequality, this implies that $|R|=|D|$ and ${}^zR=\Delta D$. Thus, $\Delta D$ is a vertex of $M$.

Finally, the indecomposable $k[\Delta D]$-module $U$ is a direct summand of $\Res_{\Delta D}^{G\times H}(M)$ and has vertex $\Delta D$. Therefore, $U$ is a source of
$M$. In other words, via the identification
$D\cong\Delta D$, the cap of $\Res_D^{ND}(V)$ is a source of the bimodule inducing the Morita equivalence.
\end{proof}

